% Figaro Paper E: On-Chain Evidence, Off-Chain Adjudication
% A Layered Approach to Smart-Contract Dispute Resolution
% Audience: legal scholars, COALA-adjacent practitioners, on-chain
% dispute-resolution researchers.

\documentclass[11pt,a4paper]{article}

\usepackage[utf8]{inputenc}
\usepackage[T1]{fontenc}
\usepackage{amsmath,amssymb,amsthm}
\usepackage{mathtools}
\usepackage{booktabs}
\usepackage{graphicx}
\usepackage{hyperref}
\usepackage{cleveref}
\usepackage{enumitem}
\usepackage{xcolor}
\usepackage[margin=1in]{geometry}
\usepackage{natbib}
\bibliographystyle{plainnat}

\theoremstyle{remark}
\newtheorem{remark}{Remark}[section]

\newcommand{\figaro}{\textsc{Figaro}}

\title{%
  On-Chain Evidence, Off-Chain Adjudication:\\
  A Layered Approach to Smart-Contract Dispute Resolution%
}

\author{
  Alessandro Daliana\thanks{Corresponding author. \texttt{adaliana@figaro.org}}
}

\date{April 2026}

\begin{document}

\maketitle

% ════════════════════════════════════════════════════════════════════
\begin{abstract}
This paper makes a focused legal argument about smart-contract
dispute resolution. The dominant on-chain dispute-resolution
designs---Kleros's Schelling-jury arbitration, Aragon Court's
DAO-governed jury, optimistic-rollup challenge windows---attempt
to perform adjudication on chain. We argue that this is a
category error: settlement is on-chain by nature, adjudication is
off-chain by nature, and conflating the two re-introduces the
discretionary intermediary that smart contracts were supposed to
remove.

We develop the argument by reference to a recently developed
settlement primitive (mechanism paper), formally
verified in the implementation paper, whose three-layer enforcement
architecture explicitly separates the bonding game (Layer 1), the
peer-coordination subgame (Layer 2), and the evidentiary record
that supports off-chain proceedings when the first two layers fail
(Layer 3). The contribution of this paper is to develop Layer 3 in
detail: the evidentiary properties of the on-chain record, the
relationship between those properties and existing
electronic-evidence law (eIDAS, UCC §1-201, the Singapore
Electronic Transactions Act, the UNCITRAL Model Law on Electronic
Commerce), and the jurisdictional implications of decoupling
settlement from adjudication.

The argument is not that smart contracts replace courts. It is
that smart contracts can produce a form of evidence that courts
already know how to use, and that the cleanest division of labor
is one in which the protocol does what it does well (deterministic
settlement) and the legal system does what it does well
(contextual adjudication). The protocol is silent on what counts
as ``satisfactory performance'' or ``material breach''---those are
adjudicative determinations---but loud on the facts that
adjudication needs as input.
\end{abstract}

\medskip
\noindent\textbf{Keywords:}
electronic evidence, smart-contract law, dispute resolution,
blockchain evidence, eIDAS, UNCITRAL, jurisdictional flexibility

% ════════════════════════════════════════════════════════════════════
\section{Introduction}
\label{sec:intro}

Smart-contract dispute resolution is currently a contested design
space. The dominant architectures attempt to perform adjudication
on chain, typically through some variant of token-weighted juror
selection or DAO governance. Kleros~\citep{kleros2019} uses
Schelling-point juror voting; Aragon Court~\citep{aragon2017}
embeds arbitration in DAO infrastructure; optimistic
mechanisms~\citep{optimism2021,uma2020} use challenge windows in
which assertions stand unless contested. Each of these
approaches addresses the dispute-resolution problem by importing
the discretionary actor that would otherwise reside in a court:
the juror, the DAO voter, the challenger.

This paper argues that the import is a mistake. The claim is not
that on-chain juries are bad mechanisms (they are interesting
mechanisms with documented properties); it is that they are
solving the wrong problem. The right problem is not ``how do we
adjudicate disputes on chain'' but ``how do we produce evidence
that off-chain forums can use to adjudicate efficiently.'' The
former requires re-creating, in adversarial token-economic
settings, the institutional infrastructure (impartial juror
selection, jurisdictional fit, due process) that civil legal
systems have refined over centuries. The latter requires only that
the on-chain artifact meet existing evidentiary standards---a
substantially easier problem, and one for which the legal
infrastructure already exists.

We make the argument in the context of \figaro{}, a settlement
primitive whose mechanism is established
in the mechanism paper and whose implementation is
verified in the implementation paper. \figaro{} explicitly does not
attempt on-chain dispute resolution. Its three-layer enforcement
architecture treats Layer 1 (bonding) and Layer 2 (peer
coordination) as the primary defenses, and Layer 3 (immutable
evidence) as the residual mechanism for cases where Layers 1 and
2 are insufficient. The present paper develops Layer 3.

\paragraph{What this paper is not.}
This paper does not provide legal advice. It does not predict the
admissibility of any specific on-chain artifact in any specific
jurisdiction, nor does it claim that the on-chain record is
itself sufficient for any particular legal claim. It argues for a
design discipline (do not perform adjudication on chain) and
catalogues the evidentiary properties of the artifact that the
discipline produces. Practitioners who wish to rely on those
properties in specific proceedings should obtain jurisdiction-specific
counsel.

\paragraph{Paper organization.}
\Cref{sec:layering} states the layered architecture.
\Cref{sec:evidentiary} develops the evidentiary properties of the
on-chain record. \Cref{sec:legal-fit} compares those properties to
existing electronic-evidence law.
\Cref{sec:jurisdiction} addresses jurisdictional implications.
\Cref{sec:against-on-chain} argues against on-chain adjudication
as a category. \Cref{sec:conclusion} concludes.

% ════════════════════════════════════════════════════════════════════
\section{The Three-Layer Architecture}
\label{sec:layering}

\figaro{}'s enforcement architecture is layered. We summarize it
here from the legal-design perspective; the mechanism details are
in the mechanism and implementation papers.

\paragraph{Layer 1 --- bonding.}
Each party to a transaction locks $2\times$ the transaction value
(asymmetric for $N$-party processes; the seller's bond scales with
upstream value). The locked capital is recoverable only through
buyer-initiated atomic resolution. Cooperation is the strictly
dominant strategy; defection costs the full bond. The vast
majority of disputes are prevented at this layer because the
incentive to defect is structurally negative.

\paragraph{Layer 2 --- peer coordination.}
Atomic resolution induces a weakest-link subgame among
co-sellers in a process tree: any one seller's defection causes
all sellers to lose their bonds. Honest sellers therefore have
direct economic incentive to monitor and pressure underperforming
co-sellers. This layer handles disputes that survive Layer 1 by
internalizing them as intra-cohort coordination problems.

\paragraph{Layer 3 --- evidence.}
The on-chain record. Every commitment, every bond deposit, every
resolution (or its absence) is recorded with block timestamps and
cryptographic signatures. When Layers 1 and 2 fail---when a party
acts irrationally, when a co-seller cohort cannot self-coordinate,
when external circumstances (regulatory action, force majeure,
participant death) prevent normal resolution---Layer 3 produces
the artifact that off-chain dispute forums need to adjudicate.

\paragraph{What is and is not on chain.}
Settlement is on chain. Adjudication is off chain. The split is
deliberate. Settlement is deterministic, automated, and
public---all properties that on-chain execution provides natively.
Adjudication is contextual, discretionary, and
jurisdiction-specific---all properties that on-chain execution is
poorly suited for. The split lets each layer do what it does best
and avoids forcing one to imitate the other.

% ════════════════════════════════════════════════════════════════════
\section{The Three Layers in Operation}
\label{sec:three-layers-operation}

The previous section states the three-layer architecture
abstractly. We walk it through operationally here, both because
the layered architecture is the paper's central thesis and because
practitioners frequently arrive with the question ``where does my
dispute go?'' that the abstract treatment does not directly answer.

\subsection{What disagreements look like at Layer 1}

Layer 1 disagreements are the ones the bonding rule prevents from
happening at all. A buyer who has signed a commitment for a
payment $P$ and posted a $2P$ bond faces a choice between
cooperation (recover the bond minus the net flow $P$) and
defection (forfeit the bond entirely). The cooperation/defection
gap is large in absolute terms ($P$ for the buyer; $P + 2G$ for
the seller, where $G$ is the cumulative upstream value); the gap
is tight in the technical sense (cooperation strictly dominates
defection when the counterparty cooperates, and is equivalent
when the counterparty defects). Most participants do not defect,
because doing so is locally irrational under the bonded posture.

\emph{What does and does not reach a forum at Layer 1.} The
clearest case for Layer-1 prevention is the
\emph{cooperation/defection} dispute: a counterparty who would
have walked away from a conventional contract because the cost of
performance exceeded the cost of breach is, under the bonded
posture, prevented from doing so because breach now costs the
forfeited bond. This class of dispute is what the bonding rule
forecloses, and it is plausibly the largest class in volume.
Real commerce, however, is dominated by a different class:
\emph{information-asymmetry} disputes about whether the seller's
performance conformed to the agreement (did the goods conform?
was time of the essence? did the description match?). The bonded
posture pressures the seller toward addressing dissatisfaction,
because the seller's bond remains locked while the buyer
withholds resolution; but it does not by itself resolve the
underlying factual question of whether the buyer's dissatisfaction
is justified. Information-asymmetry disputes survive Layer~1 by
construction and flow to Layer~2 if multi-party, or to Layer~3
if the parties cannot converge. The buyer-dominance option (do
not call \texttt{resolveProcess} until dissatisfaction is
addressed) is the principal bonded-architecture mechanism for
Layer-1 mediation of information-asymmetry disputes; whether it
suffices is a function of the dispute's specifics.

\subsection{What disagreements look like at Layer 2}

Layer 2 disagreements are multi-party coordination disputes that
the weakest-link subgame internalizes among the seller cohort.
Consider a process tree with a buyer, a primary seller, and three
sub-sellers, where one sub-seller is performing poorly enough to
threaten the parent process's resolution. Each of the three other
sellers (the primary seller and two competent sub-sellers) faces
a bond loss of $P_i + 2G_i$ if the underperformer's performance
causes the buyer to refuse resolution. They have
direct economic incentive to apply pressure on the underperformer:
phone calls, replacement-bond financing, escalation through the
cohort's coordination channels.

\emph{What this looks like operationally.} The pressure is
endogenous to the bonded structure; no governance, no platform,
no aggregator, no foundation organizes it. The companion
mechanism paper develops the result formally as the
\emph{Endogenous Coordination Pressure} proposition; from a
practical standpoint, the cohort's behavior is the operational
form of the result.

\emph{Replacement-bond patterns under exploration.} Several
runtime-tier composition patterns have been proposed for
operationalizing Layer~2 mediation but should be understood as
exploratory rather than as established mechanisms in the
mechanism-paper sense. \emph{Seller substitution} (the buyer
commits a replacement order with a new seller before
resolution; the underperforming seller's bond forfeits at
non-resolution) is the cleanest pattern, since each commitment
remains bilateral and bonded. \emph{Assembly-aggregated bonding}
(the cohort posts a shared bond pool against cohort performance)
shifts coordination from per-seller to per-cohort and raises a
distinct set of incentive-compatibility questions the
mechanism-paper analysis does not directly address.
\emph{Sponsor-of-bond patterns} (a financially-strong party
posts bond on behalf of an under-capitalized cohort member)
shade into unbonded-actor territory the no-escape-hatches
discipline rules out unless the sponsor's own bond is
posted alongside, and the sponsor relationship is itself a
bonded commitment. Each pattern is a runtime-tier composition
above the kernel; the present paper notes them as patterns
under exploration and defers full incentive analysis to the
companion runtime-composition paper (Paper~N).

\subsection{What disagreements look like at Layer 3}

Layer 3 disagreements are the residual cases that survive
Layers 1 and 2: a participant acts irrationally, the cohort cannot
self-coordinate, external circumstances prevent normal resolution,
or the underlying agreement is contested on substantive grounds
(duress, mistake, frustration, illegality, public policy). The
participants take the bonded record to an off-chain forum and the
forum adjudicates on the evidence.

\emph{The expected rate.} The Grameen joint-liability microfinance
literature reports headline default-rate reductions on the order
of magnitude of $20\% \to 2\%$ under joint-liability
arrangements with peer monitoring \citep{grameen1999};
\citet{morduch_1999} contests the headline figures on accounting
grounds (loan rescheduling, definition-of-default ambiguity,
selection effects). We do not treat the $2\%$ figure as a
peer-reviewed empirical estimate; we treat the
$20\% \to 2\%$ comparison as a stylized indication of the order
of magnitude joint-liability peer pressure has been claimed to
achieve, recognizing that the headline figure is contested. The
bonded primitive's peer-monitoring effect is structurally
stronger than Grameen's on three margins (the externality
magnitude scales with chain depth; the punishment technology
is mechanical rather than social; information availability is
on chain), and the on-chain default-and-resolution event log is
not subject to the accounting ambiguities Morduch identifies in
Grameen's self-reporting. The companion behavioral-game-theory
paper develops the benchmark in detail and is the right
reference for readers who want to push the empirical question
further. We do not commit to a specific rate prediction.

\emph{The forum the parties chose.} The forum is named by the
parties' off-chain agreement. Layer 3 does not impose a forum;
it produces the evidentiary record on which the chosen forum
operates. \Cref{sec:dispute-systems} below catalogues the
forums most commonly composed with the protocol.

\subsection{The walkthrough as practical guidance}

The three-layer architecture's practical guidance is therefore:
expect Layer 1 to handle the vast majority of cases; expect
Layer 2 to handle the multi-party residual; expect Layer 3 to
handle the small residual that survives both. Plan
infrastructure accordingly: most operational deployment work goes
into the runtime tier's Layer-1 and Layer-2 surfaces (cart UX,
attestation flows, replacement-bond composition, dispute-staging
runtime affordances), with Layer 3 being structurally reachable
but rarely invoked. A deployment that finds itself routinely
escalating disputes to Layer 3 is typically signaling that the
upstream layers are mis-configured (insufficient bonding, broken
attestation flows, mis-specified clauses), not that the
architecture is failing in the way conventional dispute-rich
commerce systems would manifest.

% ════════════════════════════════════════════════════════════════════
\section{Evidentiary Properties of the On-Chain Record}
\label{sec:evidentiary}

Before enumerating the evidentiary properties of the on-chain
record, we draw a distinction the remainder of this section
relies on. The ``on-chain record'' is not a single artifact
class; it is two artifact classes with substantially different
evidentiary character.

\paragraph{Class A: settlement byproduct.}
Commitment payloads, bond deposits, and resolution calls. These
are produced as a mechanical side effect of the bonding game: no
party can collect the settlement payout without the corresponding
on-chain events having been emitted. Their integrity as evidence
is underwritten by the game itself --- an actor cannot opt out of
producing them without opting out of settlement.

\paragraph{Class B: discretionary attestation.}
Schema-typed lifecycle events, disclosure records, proximity
witnesses, and other party-generated attestations emitted by the
attestation coordinator. These are voluntary: a party chooses
whether to attest, when to attest, and what content to place in
the attestation. The schema validator ensures the content is
well-formed; it does not and cannot ensure the content is
accurate. A party who selectively attests (silence about
underperformance, precise attestation of performance) can produce
a Class B record that survives authentication but misrepresents
reality.

\medskip

The six properties enumerated below apply \emph{unconditionally}
to Class A and \emph{conditionally} to Class B: each Class B
attestation carries the integrity of the signature and timestamp
(Class A characteristics), but the accuracy of its content
depends on the attestor's incentive structure at the moment of
attestation. A cryptoeconomics reader will recognize Class B as
subject to the same self-reporting incentive problems that
attach to every voluntary disclosure; adjudication should treat
Class B under ordinary authentication-and-hearsay doctrine
rather than extending Class A's byproduct assurance to it. An
attestation-bond layer (stake posted per attestation, slashable
on proven falsity) is a plausible Class B strengthening but lies
outside the kernel and is not yet developed in the protocol.

\paragraph{Evidentiary completeness and constitutive primacy.}
A further distinction matters for certain downstream uses:
the on-chain record is not merely \emph{admissible} as
evidence but \emph{constitutive} of the accounting record of
the process. The companion accounting
paper develops the argument that
a \figaro{} process is a self-closing ledger period in the
precise accounting sense --- the custodial balance of the
contract is the balance-sheet position; commit and resolve
events are journal and closing entries; the process closes
when resolution discharges all custodial liabilities. For
disputes whose subject matter is the accounting record of
economic activity settled through the protocol, the on-chain
record is not corroborating evidence for a party-held book;
the chain is the book, and any party-held document is a
summary or translation of it. The admissibility analysis of
the present section applies to the record's evidentiary
weight in proceedings; the constitutive analysis applies to
the record's primacy within the economic activity it
documents.

The on-chain record produced by \figaro{} has six properties that
together constitute its evidentiary value.

\paragraph{(1) Timestamped.}
Every state transition is associated with a block timestamp. The
timestamp is set by block proposers under consensus rules; it is
not subject to retroactive modification. For practical purposes,
the timestamp on a finalized block is a qualified time stamp in
the eIDAS sense.

\paragraph{(2) Cryptographically signed.}
Every commitment is signed by both parties using EIP-712 typed
structured data~\citep{eip712}. The signatures bind the commitment
content (payment, parties, denomination, agreement hash, deadline)
to the signers' private keys. Recovery is deterministic; tampering
with any field invalidates the signature.

\paragraph{(3) Content-addressed.}
Order identifiers are content hashes
($\text{orderHash} = \mathrm{keccak256}(\text{processId} \;\|\;
\text{structHash})$). The same commitment content always produces
the same identifier; different content always produces different
identifiers (collision-resistant under the assumed hash
function). This makes the record unforgeable in a specific sense:
no second commitment can claim the same identifier.

\paragraph{(4) Immutable.}
On a finalized block, the record cannot be modified or deleted. A
party seeking to repudiate a commitment cannot purge the chain of
the relevant log entries. This is a stronger property than what
electronic records normally enjoy.

\paragraph{(5) Publicly verifiable.}
Anyone can read the chain and replay the verification logic. The
evidence is not held by any single party; it is held by the
network. Disputing parties therefore cannot disagree about what
the record says---they can disagree only about what the record
\emph{means}, which is precisely the question adjudication is
suited to.

\paragraph{(6) Contains commitment intent.}
The signed payload includes a content hash $h$ of the off-chain
agreement. The off-chain agreement itself can be anything the
parties choose---a structured contract, a free-text PDF, a
schema-typed JSON document. The on-chain record does not interpret
the agreement; it commits both parties to a specific document via
the hash. A dispute over performance is then a dispute over the
referenced document, with the on-chain hash establishing
\emph{which} document is at issue.

These six properties, combined, produce an artifact that is
substantially \emph{stronger} than the electronic records most
courts already accept. Courts admit emails, electronic signatures,
and database extracts as evidence subject to authentication
challenges; the on-chain record satisfies many authentication
requirements at the cryptographic level rather than at the
custodial level.

% ════════════════════════════════════════════════════════════════════
\section{Class A and Class B Records: A Taxonomy}
\label{sec:taxonomy}

The on-chain record's evidentiary use turns on a distinction that
we have used inline in earlier sections and that deserves explicit
treatment. We taxonomize the protocol's records into two classes
distinguished by how they are produced and what they prove.

\paragraph{Class A: kernel byproducts.}
Class~A records are emitted by the kernel itself as a byproduct of
\texttt{commit} and \texttt{resolveProcess}. They include the
commitment digest (the EIP-712 hash of the typed agreement
struct), the bond deposits, the cumulative-value accumulator
state, the order status transitions, and the resolution event
with its complete payout vector. Class~A records share four
properties.

\emph{Discretion-free.} The kernel emits Class~A records without
any participant choice; if the operation succeeded, the record is
emitted in the form the kernel's transition function specifies.
There is no off-chain choice about whether to record, what to
record, or how to format the record.

\emph{Cryptographically authenticated.} Each Class~A record is
signed by the parties who participated in the operation that
produced it (EIP-712 signatures on commits; the buyer's
transaction signature on resolutions). The cryptographic
authentication is at the kernel level; an off-chain forum can
verify the signatures without consulting any oracle or
attestation service.

\emph{Append-only.} Class~A records cannot be modified after
emission. The on-chain ledger's immutability properties apply
fully; a Class~A record at block $b$ is the same record at block
$b + N$ for any $N$.

\emph{Schema-typed at the kernel level.} The structure of a
Class~A record is determined by the kernel's typed events; the
fields are known in advance and machine-readable.

The four properties together make Class~A records \emph{strong
evidence}: an off-chain forum can rely on the record's existence
and content with high confidence and minimal expert testimony.
The MLETR functional-equivalence analysis (\Cref{sec:legal-fit})
applies directly to Class~A records, with the ``reliable method''
test discharged by the kernel's bytecode and the consensus rules
of the underlying chain.

\paragraph{Class B: discretionary attestations.}
Class~B records are produced by participants through the
\texttt{AttestationCoordinator} contract under
schema-validation discipline. They include delivery
confirmations, GHG disclosures, handoff attestations,
custody-of-seal events, lifecycle stage transitions, and the
other schema-typed records that operational deployments produce.
Class~B records share four properties that contrast with Class~A.

\emph{Discretionary.} A participant chooses whether to file a
Class~B attestation, when to file it, and (within the schema's
content space) what specifically to attest. The kernel does not
require any specific Class~B record to be filed for any specific
process; whether a deployment requires a particular Class~B
record is determined by the off-chain agreement and the schema
clauses committed at signing time.

\emph{Cryptographically authenticated, plus validator-gated.}
Class~B records carry the participant's EIP-712 signature
(authenticating who filed) plus the schema-validator's gate
(authenticating that the content is well-formed under the
registered schema). The combination produces stronger
authentication than the participant's signature alone, but the
content itself is participant-asserted and may be contestable on
substantive grounds (was the delivery actually made? did the GHG
disclosure reflect actual emissions?). The companion paper on
schema design (Paper~N) develops the validator-gate's role in
the protocol-tier integrity surface.

\emph{Append-only.} Class~B records, like Class~A, cannot be
modified after emission. Participants who attest incorrectly
cannot retract; they can only file additional attestations
(typically under a different schema or stage) that contextualize
or supersede the earlier one.

\emph{Schema-typed at the protocol level.} The structure of a
Class~B record is determined by the registered schema's
specification (Layer~A in the four-layer verification stack),
ABI-encoded for on-chain validation (Layer~C), and rendered for
off-chain inspection through the runtime tier's schema-aware
modules.

The four properties together make Class~B records \emph{useful
but contestable evidence}: an off-chain forum can rely on the
record's existence (cryptographic), its participant
(cryptographic), and its formal correctness under the schema
(validator-gated), but the substantive accuracy of what the
record asserts (the delivery did happen as the seller stated;
the GHG disclosure does reflect reality) raises questions the
audit-evidence framework of \citet{aicpa2020auc500} treats
across the assertion categories of existence (did the recorded
event happen?) and accuracy (was it recorded at the right
amount, timing, and characterization?). For Class~B records the
existence and accuracy assertions are both at stake; the wider
seven-assertion AU-C~500 framework (existence, completeness,
accuracy/valuation, cutoff, classification, occurrence,
rights/obligations) supplies the language an audit-aware reader
will reach for. We use ``existence-and-accuracy'' as shorthand
for that broader frame.

\paragraph{The MLETR fit by class.}
The MLETR analysis applies most directly to Class~A records,
which carry the kernel's structural properties (singularity
through the cumulative-value accumulator; control through
buyer-dominance; integrity through the chain's consensus) and
discharge MLETR Articles~10--11's control and singularity tests
in their own right. Class~B records do not generally meet the
MLETR test on their own, because they are not transferable
records in the Article-10 sense: a delivery confirmation, a GHG
disclosure, or a custody-of-seal event is not an instrument of
title that participants treat as functionally equivalent to a
negotiable paper. Class~B records are evidentiary artifacts
\emph{about} the underlying transferable record (the
commitment), not separate transferable records inheriting the
Article 10--11 status. Their evidentiary weight in any forum
that imports MLETR-style functional-equivalence reasoning is as
corroboration to the Class~A commitment, conditional on the
schema-validation surface's correctness (a Class~B record under
a schemaId without a registered validator does not get filed in
the first place; the validator-gate at the
\texttt{AttestationCoordinator} reverts). The schema-design
discipline of Paper~N is therefore part of the corroborative
weight Class~B records carry, not part of an MLETR
transferable-record analysis applied to them directly.

\paragraph{Why the taxonomy matters for forum strategy.}
A litigant taking the bonded record to a forum should expect the
forum's treatment of Class~A and Class~B records to differ in
detail even if both are admissible.

For Class~A, the forum's analysis typically reduces to:
authentication (signatures verify; chain matches; expert if
needed), then content interpretation (what does the recorded
state mean for the parties' claims). The authentication step is
nearly mechanical; the content-interpretation step is the work.

For Class~B, the forum's analysis typically expands to include:
authentication (as for Class~A), schema-validity (was the record
filed under a registered schemaId; did it pass the validator),
then content interpretation \emph{plus substantive truth-of-content}
(does the attestation reflect the off-chain reality it asserts).
The substantive truth-of-content question engages the AU-C~500
existence and accuracy assertions at the per-attestation
granularity.

The difference matters operationally because a deployment that
relies on Class~B records to make load-bearing factual claims
(e.g., a delivery attestation that determines whether the
buyer's resolution is justified) is exposed to a different
litigation profile than one that relies on Class~A records (e.g.,
the resolution event itself, which is undeniable on chain). Most
deployments use both classes; understanding the litigation
profile of each is part of the runtime tier's design discipline.

% ════════════════════════════════════════════════════════════════════
\section{Legal Fit: Existing Electronic-Evidence Frameworks}
\label{sec:legal-fit}

We sketch how the six properties above interact with four major
electronic-evidence frameworks. The treatment is illustrative; a
full doctrinal analysis is jurisdiction-by-jurisdiction.

\paragraph{eIDAS (EU Regulation 910/2014).}
eIDAS distinguishes electronic signatures, advanced electronic
signatures (AdES), and qualified electronic signatures (QES).
EIP-712 signatures meet the AdES technical requirements (uniquely
linked to the signatory; capable of identifying the signatory;
created using means under the signatory's sole control;
linked to the signed data such that any subsequent change is
detectable). They do not, by themselves, meet the QES
requirements (which include a qualified certificate from a
trust service provider). However, the AdES status is sufficient
for most contractual purposes, and the eIDAS framework explicitly
provides that the legal effect of an electronic signature shall
not be denied solely because it is in electronic form. eIDAS
also defines qualified electronic time stamps; block timestamps
on chains with public consensus rules satisfy the technical
requirements but typically not the trust-service-provider
designation. The practical effect: the on-chain record is
admissible electronic evidence under eIDAS, with weight
determined by the court's contextual assessment.

\paragraph{UCC §1-201 and the Federal Rules of Evidence (US).}
The Uniform Commercial Code defines ``record'' as ``information
that is inscribed on a tangible medium or that is stored in an
electronic or other medium and is retrievable in perceivable
form.'' On-chain log entries satisfy this definition without
ambiguity. UCC §1-201 also defines ``signature'' inclusively
(``any symbol executed or adopted with present intention to
adopt or accept a writing''); EIP-712 signatures, being
intentionally executed by the signer's private key on the typed
commitment, qualify. The Uniform Electronic Transactions Act
(UETA) and the federal E-SIGN Act extend the same principle to
contracts more broadly.

For evidentiary admission specifically, the operative provisions
are Federal Rules of Evidence~901(b)(9), 902(13), and 902(14),
together with the \emph{Daubert} standard for any expert
testimony required to authenticate the cryptographic process.
\emph{FRE~901(b)(9)} permits authentication of evidence by
showing that ``a process or system'' produces an accurate
result; on-chain records authenticate under this rule on the
strength of the consensus protocol that produced them combined
with public availability of the bytecode that processed the
inputs. \emph{FRE~902(13)} (added 2017) permits self-authentication
of records ``generated by an electronic process or system,'' and
\emph{FRE~902(14)} permits self-authentication of data copied
from electronic devices through a process producing a digital
identifier; both rules were drafted with electronic forensic
evidence in mind and apply directly to on-chain records, where
the ``digital identifier'' is the cryptographic hash of the
record itself \citep{frelegislative2017}. The combined effect
is that the on-chain record can be self-authenticated under
902(13)--(14) in many cases, with 901(b)(9) and Daubert
available as the cross-examination route where an opposing
party challenges the authentication. The practical effect: on-chain
commitments are records and signatures under US commercial law,
admissible under the FRE through the established
electronic-evidence pathway, and authenticatable through
\emph{Daubert}-qualified expert testimony where authentication
is challenged.

\paragraph{Singapore Electronic Transactions Act (Cap. 88) and MLETR.}
Singapore's ETA closely follows the UNCITRAL Model Law and
provides functional equivalents for ``writing'',
``signature'', and ``original'' in electronic form. The ETA was
amended in 2021 to recognize electronic transferable records,
implementing the UNCITRAL Model Law on Electronic Transferable
Records (MLETR, 2017) and extending the ETA's scope to
instruments of title. On-chain records satisfy the writing and
signature requirements directly; the transferable-record
provisions warrant closer treatment because the doctrinal
question they raise is the one most directly relevant to the
process-tree structure of \figaro{}.

MLETR Article~10 requires functional equivalence with a paper-based
transferable document on three dimensions: (i) the information
contained in the electronic record is accessible for subsequent
reference; (ii) the record uses a reliable method (a) to identify
the record as the electronic transferable record, (b) to render
the record capable of being subject to control from its creation
until it ceases to have effect, and (c) to retain the integrity
of the record. Article~11 requires that ``control'' over the
record be functionally equivalent to possession of a paper
instrument: a single person identifiable as the controller, with
exclusive capacity to transfer that control, and the transfer
itself recorded reliably. The control and singularity tests are
the load-bearing features.

\figaro{}'s on-chain commitment satisfies MLETR Articles~10--11 in
the relevant senses. The orderHash uniquely identifies the
transferable record (Article~10(2)(a)). The bond mechanics give
the bonded party exclusive capacity to act on the record until
resolution (Article~11(1) singularity); the cumulative-value
accumulator ensures that the record cannot be duplicated or
forked without the chain rejecting the duplicate (Article~10(2)(b)
non-duplication). The integrity of the record is enforced
cryptographically through the chain's consensus and through
EIP-712 signature verification (Article~10(2)(c)). Where the
right to receive performance moves along a process tree, each
sub-order is itself a transferable record under the same
analysis, and the on-chain settlement records the transfer with
the reliable method MLETR contemplates. The
\citet{goode_gullifer_2017} treatment of transferable records
gives the standard doctrinal framework against which this analysis
is run; the on-chain settlement satisfies the framework's tests
without controversy.

For \emph{Class-A} artifacts (those produced as kernel byproducts:
commitment digests, settlement events, the cumulative-value
accumulator) the MLETR fit is direct. For \emph{Class-B} artifacts
(discretionary attestations, schema-validated content)
additional functional-equivalence analysis applies, since the
MLETR analysis above turned on the kernel's structural properties
rather than on attestation semantics; we do not extend the
analysis to Class-B here.

\paragraph{UNCITRAL Model Law on Electronic Commerce.}
Adopted in 1996 and supplemented by the 2017 Model Law on
Electronic Transferable Records, the UNCITRAL framework
articulates the functional-equivalence principle that underlies
most modern electronic-evidence regimes worldwide. The Model
Law's three core requirements---accessibility for subsequent
reference, integrity of the information, and identifiability of
the originator---are all satisfied by on-chain commitments. The
practical effect: jurisdictions that have adopted UNCITRAL-derived
electronic-transactions law (over 80 as of 2024) provide an
established framework for treating on-chain commitments as
admissible electronic evidence.

\paragraph{The summary picture.}
The on-chain record is admissible electronic evidence in every
jurisdiction whose electronic-evidence law derives from the
UNCITRAL framework or implements equivalent functional principles.
What varies across jurisdictions is the weight courts assign to
the evidence (a question of judicial practice rather than
admissibility) and the further procedural requirements for
authenticating the chain itself (typically satisfied by expert
testimony or by the court taking judicial notice of well-known
chains). No jurisdiction we are aware of categorically excludes
on-chain records from evidentiary consideration.

% ════════════════════════════════════════════════════════════════════
\section{Substantive Defenses and Post-Settlement Restitution}
\label{sec:defeasance}

A reader sympathetic to the argument so far will raise the
sharpest question this paper has to answer. Suppose the on-chain
settlement is unimpeachable as evidence; suppose every framework
above admits it. What happens when an external court, applying
the substantive contract law of its jurisdiction, finds that the
underlying agreement was procured under duress, formed under
fundamental mistake, frustrated by impossibility, voided by
illegality, or unconscionable in its terms? The on-chain
settlement has, by hypothesis, already executed: bonds are
released, payouts distributed, and the kernel's state shows the
process resolved. The question is what the court can do with
that fact and whether the substitution between settlement
(on-chain) and adjudication (off-chain) the present paper has
defended is consistent with the substantive-defense
jurisprudence the dominant tradition has developed over centuries.

We treat the question in three registers: the doctrinal
inventory, the substantive-defenses analysis, and the
restitutionary mechanics by which a court order to unwind
interacts with an already-discharged kernel state.

\paragraph{Doctrinal inventory.}
The defenses that void or rescind a contract \emph{ab initio},
or that excuse subsequent performance, are doctrinally settled
in the major commercial jurisdictions. Common-law systems
recognize: \emph{duress} (consent obtained by improper threat);
\emph{undue influence} (consent obtained through abuse of a
relationship of trust); \emph{misrepresentation} (consent
obtained by false statements about material facts);
\emph{fundamental mistake} (consent given under a shared mistake
about a fact essential to the agreement);
\emph{unconscionability} (terms so one-sided that enforcement
shocks the conscience of the court);
\emph{frustration of purpose} (subsequent events render
performance fundamentally different from what was undertaken);
\emph{impossibility / impracticability} (subsequent events
render performance physically or practically impossible);
\emph{illegality} (the agreement requires conduct prohibited by
law); \emph{public policy} (enforcement would violate public
policy independent of any specific statute). Civil-law systems
recognize closely parallel doctrines under different names
(\emph{vices du consentement}, \emph{erreur},
\emph{circonstances exceptionnelles}, etc.). The doctrines are
substantive: they apply to \emph{the agreement} the bonded
commitment records, not to the bonded commitment as a
technological artifact.

\paragraph{Substantive-defenses analysis.}
The bonded primitive is silent on every doctrine in the inventory.
The kernel does not know whether consent was procured under
duress, whether parties were mistaken about a material fact,
whether the underlying transaction was illegal, or whether
performance has been frustrated by external circumstance. The
kernel records what was signed; it makes no representation about
the conditions under which signing occurred or the legality of
what was signed for. This is by design --- the kernel cannot make
substantive-defense determinations without doing the kind of
discretionary work the no-escape-hatches design principle of the
mechanism paper rules out --- and it is by design also a
\emph{feature} for the present paper's purpose, because it leaves
the substantive-defense apparatus exactly where the dominant
tradition has placed it: in courts applying their own
constitutional and doctrinal apparatus to the agreement
recorded in the bonded commitment.

The leading case-law authority on smart-contract substantive
defense is \emph{Quoine Pte Ltd v B2C2 Ltd}
\citep{quoine_b2c2_2019, quoine_b2c2_2020}. The Singapore
International Commercial Court (and on appeal the Court of
Appeal) held that a smart-contract execution can be subject to
the equitable doctrine of unilateral mistake, and that a party
who knows or ought to know of a counterparty's mistake at the
moment a smart contract executes can be required to make
restitution. The Quoine line is the authoritative answer that
substantive-defense doctrines apply to smart-contract executions
in the jurisdictions whose courts are willing to recognize them,
and the bonded primitive is no exception.

\paragraph{Restitutionary mechanics.}
The mechanics are where the question becomes operationally
sharp. Three sub-cases.

\emph{Case 1: court orders specific re-transfer off-chain.}
Where the court determines that one party should not have
received the on-chain payout (e.g., because the underlying
agreement was void for fundamental mistake), the standard remedy
is a restitutionary order requiring that party to transfer the
received value back to the other off-chain. This remedy operates
on the parties' real-world capacity to comply, not on the kernel
state. The kernel records show the original settlement
unchanged; the off-chain transfer is recorded separately,
typically through conventional financial rails or through a
fresh on-chain transaction. The kernel state is, in this
arrangement, evidentiary input to the dispute and operational
substrate for the off-chain remedy; it is not itself the locus
of the unwinding. This is the most common form of post-settlement
restitution and the form for which the existing legal apparatus
is well-equipped.

\emph{Case 2: court orders replacement-bonded counter-process.}
Where ongoing performance is required and a counter-flow is
appropriate, the parties (or the court via an order on the
parties) can constitute a fresh \figaro{} process that flows
value in the corrective direction. This uses the bonded
primitive as the implementation of the court's remedy rather
than as an obstacle to it. The original process remains
resolved; the new process is independent and bonded under the
same kernel.

\emph{Case 3: court orders an arrangement the kernel cannot
support.}
Where the court's appropriate remedy cannot be implemented
within the kernel's bilateral primitive (e.g., because it
requires the kind of discretionary administrative authority the
kernel forbids), the remedy operates on the parties' off-chain
assets and obligations through conventional mechanisms (asset
seizure, contempt sanctions, criminal liability for non-compliance
with civil orders). The kernel does not obstruct this remedy;
it does not implement it either. This is consistent with the
general principle that the substrate is sub-political, not
anti-political (\Cref{sec:against-on-chain} returns to this).

\paragraph{Implication for Paper~A's escape-hatch result.}
The mechanism paper's \emph{escape-hatch weakness} theorem
(Theorem~4.7 of Paper~A) and the carve-out in its Remark~4.8
deserve a more precise characterization than either text gives.
The theorem rules out unilateral exit paths internal to the
mechanism --- timeouts, arbitrators, governance votes ---
because those would weaken the Nash equilibrium. Remark~4.8
carves out external legal forums applying duress, frustration,
and similar substantive doctrines, on the ground that those
forums are constrained by their own institutional bond
structures and operate on the bonded commitment as evidentiary
input.

The carve-out is sound on closer analysis but for a different
reason than Remark~4.8 suggests. The relevant fact is not that
external forums are themselves bonded; it is that the
substantive-defense jurisprudence operates on the agreement
\emph{as recorded} after the kernel has settled, not on the
kernel itself, and produces remedies that reach the parties
through their conventional legal exposure rather than through
the bonded escrow. The kernel's equilibrium is preserved
because no remedy can reach into the discharged escrow; the
substantive defenses operate downstream of settlement. A
participant deciding whether to bond does so with the
understanding that the bonded commitment is enforceable as
between the parties and may, separately, be tested in court
under conventional contract doctrines on the same terms as any
other contract. Both the on-chain enforcement and the
off-chain doctrinal review apply; they are not in competition
because they reach different aspects of the same arrangement.
The bonded primitive does not displace contract law's
substantive doctrines; it produces a settlement under which
those doctrines continue to apply in the conventional way.

% ════════════════════════════════════════════════════════════════════
\section{A Worked Example: New York Forum, Cross-Border Process}
\label{sec:ny-worked-example}

A worked example sharpens the abstract analysis. Consider a
process tree with a buyer in New York, a primary seller in
Singapore, and a sub-seller in Argentina, settled in
\texttt{USDC} on Ethereum L1. The off-chain agreement specifies
New York Supreme Court as the forum and New York law as the
governing law (a common contractual arrangement for
cross-border commerce involving a US counterparty). A dispute
arises: the primary seller claims the buyer never resolved the
process despite delivery; the buyer claims the goods did not
conform to the agreed specification.

The bonded primitive's contribution to the dispute is direct.
The on-chain record establishes: the commitment hash matching
the off-chain agreement (UCC~§1-201 record; FRE~902(13)
self-authenticating); the buyer's and seller's signatures on the
EIP-712 typed commitment (UCC~§1-201 signature; AdES under
eIDAS; signature under E-SIGN/UETA); the cumulative-value
accumulator's monotonic state through every commit and resolve
operation; and the absence (or presence) of a
\texttt{resolveProcess} transaction. The seller's claim --- that
the buyer never resolved --- is determinable from the chain
state directly. The buyer's claim --- that the goods did not
conform --- is not determinable from the chain state and
requires off-chain evidence: delivery records, photographs,
laboratory analyses, and (where the process used the
\figaro{}-fulfilment schema) the schema-validated attestations
that may already be on chain.

The New York court, applying New York contract law (UCC
Article~2 for the goods, NY CPLR for the procedural framework,
NY GBL §75-a for electronic signatures), proceeds as follows.
Admissibility of the on-chain record is established under the
FRE pathway above (902(13) self-authentication; 901(b)(9)
process-or-system authentication; \emph{Daubert}-qualified
expert if needed). The substantive question --- did the goods
conform? --- is decided on the off-chain evidence under
ordinary contract doctrine. The remedy, if the court finds for
the buyer, is restitutionary: an order requiring the seller to
repay the unconformed portion of the payment. The order operates
on the seller's off-chain assets through standard enforcement
mechanisms (US--Singapore reciprocity, New York Convention, asset
attachment in either jurisdiction).

The bonded settlement does not interfere with any of this. It
provides the evidentiary record and the bilateral-enforcement
backstop that prevents the seller from walking away with the
funds before the dispute is heard --- the usual gap that
cross-border commerce has had to manage through letters of
credit, escrow agents, and intermediary banks. The bonded
primitive supplies that function more cleanly, leaving the
substantive doctrine of the New York court untouched. This is
the architecture the present paper has been describing: the
primitive replaces the enforcement gap that conventional
commerce has had to fill with intermediaries; it does not
replace the substantive doctrine that the court applies, and it
is no less well-positioned in cross-border settings than in
domestic ones because the evidentiary properties are the same
in either case.

% ════════════════════════════════════════════════════════════════════
\section{Jurisdictional Implications}
\label{sec:jurisdiction}

The decoupling of settlement from adjudication has substantive
jurisdictional consequences.

\paragraph{Forum is contractual.}
The on-chain commitment includes an off-chain agreement hash $h$.
That off-chain agreement can---and typically should---include a
forum-selection clause, a choice-of-law clause, and an
arbitration clause if the parties prefer arbitration to court.
The on-chain record proves what was agreed; the off-chain
agreement specifies where and under what law the agreement is
adjudicated. Parties who omit a forum clause default to whatever
private-international-law analysis the seised forum applies; this
is a known regime, not a novel problem.

\paragraph{Forum-shopping is bounded.}
A party seeking to forum-shop is bounded by (a) the
forum-selection clause in the off-chain agreement, (b) the seising
forum's own jurisdictional rules, and (c) the recognition-and-enforcement
treaties between the seising forum and any forum where the
counterparty's assets reside. None of these constraints is
weakened by the on-chain settlement; they all operate on the
off-chain agreement and the parties' real-world circumstances.

\paragraph{Evidence is jurisdiction-displaced.}
The integrity of the on-chain evidence does not depend on the
seising forum's recognition of any specific cryptographic
standard. The signatures verify, or they do not; the chain hashes
match, or they do not. A forum that doubts the cryptographic
analysis can engage an expert; a forum that does not is in the
position of any forum facing forensic evidence.

\paragraph{Cross-border processes.}
A process tree can involve parties in multiple jurisdictions
simultaneously. The on-chain settlement is jurisdiction-blind:
the bond mechanics work identically regardless of where the
parties reside. Off-chain disputes that arise from such processes
are handled by the same private-international-law machinery that
already handles cross-border disputes in conventional
commerce---typically with the on-chain record providing
substantially more evidentiary clarity than the parties would
have produced under conventional bookkeeping.

% ════════════════════════════════════════════════════════════════════
\section{Why Not On-Chain Adjudication}
\label{sec:against-on-chain}

We have argued the positive case for off-chain adjudication. We
now argue the negative case against on-chain adjudication.
Several patterns occupy this space; we treat each briefly.

\paragraph{Schelling-point juror selection (Kleros).}
A token-weighted juror panel votes on the disputed question; the
plurality wins; jurors who voted with the plurality are rewarded,
jurors who voted against are slashed. The mechanism is elegant
in its abstraction but has three problems for serious adjudication.
First, the juror panel is not selected for jurisdictional
fit---a Kleros juror in jurisdiction $A$ may know nothing of
jurisdiction $B$'s law that the parties agreed should govern.
Second, the juror's incentive is to vote with the perceived
plurality, not to vote for the legally correct outcome---a
Schelling point on plurality is structurally distinct from a
Schelling point on truth. Third, the slashing mechanism imports a
new dispute (was the slashing fair?) that has no terminal forum
on chain.

\paragraph{DAO-governed arbitration (Aragon Court).}
Aragon Court is structurally similar to Kleros, with the
juror-selection layer integrated with DAO governance. This
inherits all of Kleros's problems plus the additional problem
that DAO voting power is plutocratic in proportion to token
holdings; the wealthy juror has structurally more votes than the
poor juror. Whether or not one views this as a desirable
property of arbitration is a normative question; that it differs
from the conventional ``one juror, one vote'' premise of jury
trials is a doctrinal question that has not been settled.

\paragraph{Optimistic challenge windows.}
Optimistic mechanisms~\citep{optimism2021,uma2020} avoid the juror
problem by allowing assertions to stand unless challenged within a
fixed window. This works for settlement-layer state machines but
does not work for adjudication, because adjudication is precisely
the activity of assessing an assertion that has been challenged.
An optimistic ``adjudication'' is just an assertion with a
deadline---it has no mechanism for resolving the underlying
question, only for resolving the question of who blinked first.

\paragraph{The structural objection.}
All on-chain adjudication mechanisms share the same structural
flaw. The kernel mechanism of the mechanism paper
establishes that any unilateral exit from the bonded state
requires either a third-party decision (which introduces an
unbonded actor whose incentives are not constrained by the bond
structure) or a payoff modification that weakens the Nash
equilibrium. On-chain adjudication is the third-party-decision
case: the juror, the DAO voter, the challenger is the unbonded
actor whose decision determines the outcome. Whatever the
mechanism designer does to constrain that actor (token slashing,
Schelling incentives, reputation), the actor is still unbonded
and the kernel's self-enforcing property is broken.

\paragraph{The off-chain alternative.}
Off-chain adjudication by an existing legal forum has the
opposite structure. The forum's incentives are constrained by
its institutional context---judges face appellate review,
arbitrators face institutional reputation, both face professional
sanction. The forum is not unbonded in the kernel sense; it
operates inside a different bond structure (its institutional
authority). The off-chain evidence record is therefore consumed
by an actor whose incentives are already constrained by a regime
that has been refined for centuries. The protocol does not
re-implement that regime; it provides input to it.

\paragraph{Composition with Kleros as forum, not as kernel
adjudicator.} The objection above is to on-chain adjudication
\emph{at the kernel layer}. It is not an objection to Kleros, or
to any other adjudication mechanism, when the parties' agreement
designates it as the forum. The kernel of \figaro{} is
forum-agnostic: it does not adjudicate; it produces an evidence
record. If two parties sign an off-chain agreement whose
arbitration clause names Kleros as the dispute forum, the
evidence record produced by their \figaro{} commitments is
exported to Kleros and the dispute is adjudicated there. From
the kernel's perspective Kleros is the off-chain forum the
parties chose, in the same sense that the Singapore International
Arbitration Centre or the New York Supreme Court can be the
off-chain forum the parties choose. The structural critique above
is not an endorsement of some forums over others; any forum the
parties name in their agreement is their forum to choose, and
\figaro{} composes with it by exporting evidence to it.

\paragraph{Engagement with the legal-academic literature on Kleros.}
A small but real legal-academic literature has developed around
Kleros and its peers. \citet{ortolani2019blockchain} is the
principal engagement from the international-arbitration side,
raising the recognition-and-enforcement question under the New
York Convention~\citep{nyc1958}: can a Kleros decision be
enforced abroad as an arbitral award, and on what doctrinal
basis? \citet{kaal2020dispute} reviews the broader landscape of
decentralized dispute resolution. \citet{defilippi_wright_2018}
situates Kleros and similar mechanisms within the lex
cryptographia frame. \citet{rabinovich_einy_katsh_2017} place
them in the longer trajectory of online dispute resolution. The
literature is not yet large but it is real, and its general
posture is compatible with this paper's: most of the
legal-academic engagement with Kleros expresses concern about
adjudicative-quality and recognition issues, not about the
existence of the mechanism.

\paragraph{Our committed view on the New York Convention question.}
The Kleros recognition question deserves a committed answer
rather than a deferral. Our reading of Article~I of the New
York Convention --- which defines its scope to ``arbitral
awards made in the territory of a State other than the State
where the recognition and enforcement of such awards are
sought'' --- is that a Kleros decision is capable of falling
within the Convention's definition of an arbitral award where
the parties' arbitration agreement so designates, subject to
Article~V's grounds for refusing recognition (most relevantly
the public-policy ground of Article~V(2)(b) and the
``proper law of the arbitration agreement'' analysis under
Article~V(1)(a)). The recognition is not automatic --- a court
asked to enforce a Kleros decision can refuse on Article~V
grounds, and the substantive-defense doctrines of
\Cref{sec:defeasance} apply on the merits as they would for any
arbitral award. The recognition is also jurisdictionally
sensitive: Convention-implementing legislation varies, and
common-law and civil-law jurisdictions have different default
postures toward unconventional arbitration arrangements. Our
view is therefore that Kleros decisions \emph{can} be arbitral
awards under the Convention, conditional on competent drafting
of the arbitration agreement and on the seising court's
willingness to read the Convention's definition of an
``arbitral tribunal'' broadly enough to include the
juror-selection mechanism Kleros operates. The contrary view
\citep[on which see][]{ortolani2019blockchain} is defensible
and turns on the same Article~I and Article~V considerations.
We do not treat the question as settled, but we do treat it as
takable, and we have written the architecture to be neutral on
the choice of forum so that parties whose recognition theory
favors a state-court forum can choose one and parties whose
theory accommodates a Kleros forum can choose Kleros, with
\figaro{}'s evidence record exported to either.

The recognition-and-enforcement question raised by Ortolani and
others is in our reading the binding legal-doctrinal question,
and the layered architecture of \Cref{sec:layering} mitigates
it: the on-chain record is admissible electronic evidence under
the frameworks of \Cref{sec:legal-fit} regardless of how the
forum that consumes the evidence is constituted, and
forum-clause discipline allows the parties to name a forum
(Kleros, SIAC, court) whose decisions they have a clean
recognition theory for in the jurisdiction of intended
enforcement.

% ════════════════════════════════════════════════════════════════════
\section{Composition with Dispute Systems}
\label{sec:dispute-systems}

The previous section makes the negative case (why on-chain
adjudication breaks the equilibrium); a constructive treatment is
also owed to the practitioner who arrives with the question
``which forum should I name in my agreement, and how do I compose
with it.'' We catalogue the four classes of forum the protocol
composes with in practice and treat each as a composition
pattern.

\subsection{State courts (conventional civil and commercial
courts)}

Conventional state courts are the highest-recognition forum
available and the appropriate destination for disputes whose
substantive resolution requires the full apparatus of
contract law (duress, mistake, frustration, illegality, public
policy --- the doctrines treated in
\Cref{sec:defeasance}). The composition pattern is direct.

\emph{Forum clause.} The off-chain agreement names a specific
state court (e.g., New York Supreme Court, Singapore International
Commercial Court, English High Court) and a governing law (e.g.,
New York law, Singapore law, English law). The clause is a
standard contractual provision; the bonded primitive does not
modify the law of forum selection.

\emph{Evidence path.} The bonded record is admitted under the
relevant evidentiary framework (\Cref{sec:legal-fit}: FRE
901(b)(9) / 902(13)--(14) for US courts, eIDAS AdES recognition
for EU courts, Singapore ETA / MLETR for Singapore courts,
analogous frameworks elsewhere). Authentication challenges are
typically minimal because the cryptographic properties are
self-authenticating; substantive challenges go to the merits in
the ordinary way.

\emph{Recognition and enforcement.} A judgment from a recognized
state court is enforceable through the standard
private-international-law machinery: bilateral and multilateral
recognition treaties, the Hague Convention on Recognition and
Enforcement of Foreign Judgments (where in force), and the
forum's own domestic enforcement powers. Enforcement against
on-chain assets requires a court with jurisdiction over the
asset-holding party rather than over the bonded primitive itself
(which has no jurisdiction in the relevant sense; the kernel is
not a party).

\emph{When state courts are appropriate.} Disputes involving
substantive contract defenses, public-policy questions,
sanctions enforcement, criminal conduct adjacent to the bonded
commerce, regulatory enforcement, or any matter where the parties
need the recognition properties of a state court's judgment.

\subsection{International commercial arbitration centers}

International commercial arbitration centers (the Singapore
International Arbitration Centre / SIAC, the International
Chamber of Commerce / ICC, the London Court of International
Arbitration / LCIA, the Hong Kong International Arbitration
Centre / HKIAC, and analogous institutions) are the appropriate
destination for cross-border disputes whose parties prefer
arbitration to court litigation. The composition pattern is
similar to state courts but uses the New York Convention's
recognition machinery rather than bilateral judgment-recognition
treaties.

\emph{Forum clause.} The off-chain agreement names the
arbitration center, specifies the seat of arbitration, the
language of the proceeding, and the institution's procedural
rules (e.g., ``arbitration administered by SIAC under its 2025
Rules; seat in Singapore; English language; three arbitrators
appointed under Article 14'').

\emph{Evidence path.} Arbitration centers admit electronic
evidence broadly and typically have flexible authentication
standards; the bonded record satisfies the standards
straightforwardly.

\emph{Recognition and enforcement.} The New York Convention
\citep{nyc1958} is the canonical recognition framework for
arbitral awards; 173 state parties as of 2024. An award rendered
under SIAC / ICC / LCIA rules is enforceable in any New York
Convention state subject to the Article~V grounds for refusal
(public policy, due process, etc.). The bonded primitive
composes with this enforcement chain natively because the
underlying agreement is a recognized arbitration agreement
under Article~II.

\emph{When arbitration centers are appropriate.} Cross-border
disputes where the parties prefer institutional arbitration's
combination of expertise, procedural flexibility, and
international enforceability over the conventional-court
alternative; high-value disputes where the parties want
arbitrators with sector-specific expertise.

\subsection{Online dispute resolution (ODR) platforms}

ODR platforms (eBay's Resolution Center as the volume archetype,
Modria, smartsettle, and the broader online-dispute-resolution
literature \citep{rabinovich_einy_katsh_2017}) are the
appropriate destination for high-volume, low-value disputes
where the cost of conventional adjudication exceeds the
disputed amount. The composition pattern is structurally
distinct from courts and arbitration in that the ODR platform's
recognition properties are typically weaker but its operational
cost is much lower.

\emph{Forum clause.} The off-chain agreement names the ODR
platform and accepts its terms-of-service for dispute handling.
The platform's binding effect is contractual rather than
treaty-based.

\emph{Evidence path.} ODR platforms typically accept the bonded
record's evidentiary properties without specialized
authentication; the platform's internal review uses the record
as the dispositive content for most operational disputes
(delivery confirmation, GHG disclosure verification, handoff
attestation review).

\emph{Recognition and enforcement.} ODR decisions are typically
enforced by the platform itself (e.g., a refund or replacement
ordered by the platform's review process) rather than through
external state machinery. A practitioner pattern of pairing an
ODR clause with an arbitration clause that treats unappealed
ODR decisions as final and arbitration-recognized is
structurally available but doctrinally delicate: the arbitration
clause must be drafted to constitute the ODR review as the
arbitral process (or as a binding precursor with a separate
arbitral-confirmation step), and the resulting award's
recognition under the New York Convention turns on whether the
ODR proceeding satisfies Article~V(1)(b) due-process
requirements. Many ODR platforms' procedures (no live hearing,
no document discovery, opaque review) fail Article~V(1)(b) on
standard readings, so the hybrid is a careful drafting matter
rather than an off-the-shelf composition.

\emph{When ODR is appropriate.} Routine commercial disputes at
operational scale where the dispute volume rules out
court/arbitration economics; deployments where the parties
prefer fast resolution to definitive recognition.

\subsection{Schelling-point juror systems (Kleros and analogues)}

Kleros and analogous Schelling-point juror systems
\citep{kleros2019} occupy a distinctive position. They are not
state courts, not traditional arbitration centers, and not
conventional ODR platforms. They are mechanism-based forums
where a Schelling-point juror selection produces decisions
through coordination on the focal answer. The protocol composes
with them as the parties' chosen forum, but the recognition
and enforcement properties are the most contested of the four
classes.

\emph{Forum clause.} The off-chain agreement names Kleros (or
analogue), specifies the court within Kleros (the contracting
parties choose from the registered courts in the Kleros system),
and specifies the procedural rules (juror count, evidence-period
duration, appeal availability).

\emph{Evidence path.} The bonded record is exported to Kleros
through the parties' filing of the dispute; jurors evaluate the
evidence and vote.

\emph{Recognition and enforcement.} \Cref{sec:against-on-chain}
above developed the committed view: a Kleros decision is
\emph{capable} of falling within the New York Convention's
definition of an arbitral award where the parties' arbitration
agreement so designates, subject to Article~V grounds for
refusing recognition. \citet{ortolani2019blockchain} raises three
specific objections that the committed view must engage. First,
NYC Article~I requires an identifiable ``tribunal'', and Kleros
jurors are anonymized, rotating, and token-incentivized rather
than arbitrators in the institutional sense the Convention has
historically contemplated. Our committed view contests this
objection on the ground that the Convention's tribunal
definition is functional rather than institutional --- the
relevant question is whether a proceeding produces a decision
on the merits between the parties, not whether it does so
through institutionally-credentialed arbitrators. Second,
Article~V(1)(d) requires that ``the composition of the arbitral
authority'' was in accordance with the parties' agreement, and
algorithmic juror selection raises the question of whether
``composition'' has been agreed in the relevant sense. Our
committed view contests this objection by reading the parties'
acceptance of the Kleros court's procedural rules (which the
parties opt into when naming the court) as constituting the
required agreement on composition. Third, Article~V(2)(b)
permits refusal of recognition on public-policy grounds, and
many civil-law jurisdictions treat unreasoned awards as
presumptively contrary to public policy --- Kleros decisions
are vote tallies, not reasoned awards. \emph{Our committed view
concedes this objection in civil-law jurisdictions whose public
policy requires reasoned awards.} Recognition of Kleros
decisions in such jurisdictions is materially harder than the
generic analysis suggests, and parties enforcing in those
jurisdictions should pair the Kleros clause with a written
opinion supplement or choose a different forum. The committed
view is therefore: takable on Article~I and Article~V(1)(d) on
the functional readings above; not takable on Article~V(2)(b)
in civil-law public-policy jurisdictions whose reasoned-award
requirement applies.

\emph{When Kleros is appropriate.} Disputes whose parties prefer
the Schelling-point juror mechanism's combination of neutrality
(no state authority), speed (faster than conventional
arbitration), and the token-economic incentive structure for
juror participation, over the recognition properties of
conventional forums; operational deployments where the bond
posture itself produces strong enough enforcement that the
forum's recognition is needed only in a small minority of
cases. Unsuitable for deployments whose intended enforcement
jurisdictions include civil-law public-policy regimes that
require reasoned awards.

\subsection{Composition discipline across the four}

Two structural points apply across all four forum classes.

\emph{Forum-clause discipline is the parties' choice and the
parties' responsibility.} The bonded primitive does not
recommend a forum; it does not prefer state courts to arbitration
centers, ODR to Kleros, or any of the four to any other. The
forum is named by the parties' off-chain agreement, and the
parties bear the consequences of their forum choice (recognition
properties, enforcement chain, procedural-rule adoption,
operational cost). The architecture's role is to produce
evidence the chosen forum can use, not to constrain the choice.

\emph{Recognition theory is upstream of forum choice.} The
parties should name a forum whose decisions they have a clean
recognition theory for in the jurisdiction of intended
enforcement. A US-based seller and an EU-based buyer who name
SIAC arbitration with Singapore seat have a clean recognition
theory through the New York Convention. The same parties naming
an ODR platform without arbitration backing have no clean
recognition theory if the ODR decision is contested. The same
parties naming a state court whose judgments are not recognized
in the counterparty's jurisdiction have a recognition gap that
the bonded primitive cannot close.

\emph{The bonded primitive shifts the balance toward weaker
forums.} Because the bonded posture itself enforces most
disputes at Layers 1 and 2, the residual that reaches Layer 3
is small (\Cref{sec:three-layers-operation}). A smaller residual
volume can sometimes justify a weaker forum than the deployment
would have chosen under conventional commerce: a deployment that
would have insisted on SIAC arbitration in conventional commerce
might choose an ODR platform under the bonded primitive because
the residual volume reaching Layer 3 is small enough that ODR's
weaker recognition is acceptable. This shift is not the same as
saying the architecture replaces stronger forums; it is saying
the architecture changes the cost-benefit calculation of which
forum is appropriate for what residual.

% ════════════════════════════════════════════════════════════════════
\section{Conclusion}
\label{sec:conclusion}

We have argued for a layered approach to smart-contract dispute
resolution: settlement on chain, adjudication off chain, with a
deliberately designed evidentiary record that bridges the two.
The architecture inherits the strengths of each layer (chain
determinism for settlement, legal-system contextual judgment for
adjudication) without forcing either to imitate the other. It
also resolves the structural objection to on-chain adjudication:
the unbonded juror does not need to exist because the legally
constituted forum already does.

The contribution is modest but, we believe, important. The
debate over smart-contract dispute resolution has been dominated
by the assumption that on-chain settlement implies on-chain
adjudication. The decoupling we propose is a different design
posture, and it permits a substantially cleaner alignment with
existing electronic-evidence law.

\paragraph{What remains for legal scholarship.}
Three questions are out of scope for this paper but warrant
sustained legal-scholarly attention. First, the doctrinal
analysis of on-chain commitments under specific jurisdictions'
contract-formation, evidence, and electronic-signature laws is
substantial and necessarily jurisdiction-by-jurisdiction. Second,
the question of how courts should address chain-finality
assumptions (what happens when the underlying chain reorganizes
substantially after a contested transaction) requires technical-legal
collaboration that has barely begun. Third, the question of
whether and when the on-chain record should be treated as
self-authenticating (in the US Federal Rules of Evidence sense,
or its analogues in other common-law jurisdictions) is a question
of judicial practice that will develop with case law.

\paragraph{What remains for protocol design.}
On the protocol side, the principal open question is how schema-typed
attestations (implementation paper) interact with off-chain
dispute proceedings. The attestation surface allows parties to
record structured evidence (geo proofs, lifecycle stage
transitions, GHG disclosures) on chain in addition to the bare
commitment. How that structured evidence is treated by off-chain
forums---whether as supplementary fact, as binding stipulation,
or as expert testimony---is design space that has not yet been
explored.

\section*{Acknowledgments}
This argument owes a substantial debt to the Coalition of
Automated Legal Applications (COALA), whose work over many years
has established the framework within which questions of this kind
can be productively asked. I thank colleagues there and at the
Stanford CodeX Center for conversations that sharpened the
distinction between settlement and adjudication that this paper
develops.


% ════════════════════════════════════════════════════════════════════
% References
% ════════════════════════════════════════════════════════════════════

\begin{thebibliography}{99}

\bibitem[Aragon(2017)]{aragon2017}
Cuende, L. and Izquierdo, J.
\newblock Aragon Network: A Decentralized Infrastructure for Value Exchange.
\newblock Aragon White Paper, 2017.

\bibitem[De Filippi and Wright(2018)]{defilippi_wright_2018}
De Filippi, P. and Wright, A.
\newblock \emph{Blockchain and the Law: The Rule of Code}.
\newblock Harvard University Press, Cambridge, MA, 2018.

\bibitem[EIP-712(2017)]{eip712}
Bloemen, R., Logvinov, L., and Evans, J.
\newblock {EIP-712}: Ethereum typed structured data hashing and signing.
\newblock Ethereum Improvement Proposal, 2017.
\newblock \url{https://eips.ethereum.org/EIPS/eip-712}

\bibitem[Kaal(2020)]{kaal2020dispute}
Kaal, W.~A.
\newblock Decentralized Dispute Resolution.
\newblock University of St.\ Thomas (Minnesota) Legal Studies
  Research Paper No.\ 20-13, 2020.

\bibitem[Kleros(2019)]{kleros2019}
Ast, F. and Lesaege, C.
\newblock Kleros: A Protocol for a Decentralized Justice System.
\newblock Kleros White Paper, 2019.
\newblock \url{https://kleros.io/yellowpaper.pdf}

\bibitem[New York Convention(1958)]{nyc1958}
United Nations.
\newblock \emph{Convention on the Recognition and Enforcement
  of Foreign Arbitral Awards} (the New York Convention).
\newblock 330 U.N.T.S. 38, 1958.

\bibitem[Optimism(2021)]{optimism2021}
Optimism Foundation.
\newblock Optimistic Rollups.
\newblock Optimism Technical Documentation, 2021.

\bibitem[Ortolani(2019)]{ortolani2019blockchain}
Ortolani, P.
\newblock The Impact of Blockchain Technologies and Smart
  Contracts on Dispute Resolution: Arbitration and Court
  Litigation at the Crossroads.
\newblock \emph{Uniform Law Review}, 24(2):430--448, 2019.

\bibitem[Rabinovich-Einy and Katsh(2017)]{rabinovich_einy_katsh_2017}
Rabinovich-Einy, O. and Katsh, E.
\newblock \emph{Digital Justice: Technology and the Internet of
  Disputes}.
\newblock Oxford University Press, New York, 2017.

\bibitem[Quoine v B2C2 (2019)]{quoine_b2c2_2019}
Quoine Pte Ltd v B2C2 Ltd [2019] SGHC(I) 03 (Singapore
International Commercial Court).

\bibitem[Quoine v B2C2 (2020)]{quoine_b2c2_2020}
Quoine Pte Ltd v B2C2 Ltd [2020] SGCA(I) 02 (Singapore Court of
Appeal).

\bibitem[FRE Legislative History(2017)]{frelegislative2017}
US Judicial Conference, Advisory Committee on Evidence Rules.
\newblock Report on Proposed Amendments to Federal Rules of Evidence
  (902(13) and 902(14)).
\newblock Effective December~1, 2017.

\bibitem[Goode and Gullifer(2017)]{goode_gullifer_2017}
Goode, R. and Gullifer, L.
\newblock \emph{Goode and Gullifer on Legal Problems of Credit and
  Security}, 6th edition.
\newblock Sweet \& Maxwell, London, 2017.

\bibitem[Yunus(1999)]{grameen1999}
Yunus, M.
\newblock \emph{Banker to the Poor: Micro-Lending and the Battle
  Against World Poverty}.
\newblock PublicAffairs, New York, 1999.

\bibitem[Morduch(1999)]{morduch_1999}
Morduch, J.
\newblock The Microfinance Promise.
\newblock \emph{Journal of Economic Literature},
  37(4):1569--1614, 1999.

\bibitem[AICPA(2020)]{aicpa2020auc500}
American Institute of Certified Public Accountants.
\newblock \emph{AU-C Section 500: Audit Evidence}.
\newblock AICPA, Durham, NC, 2020.

\bibitem[UMA(2020)]{uma2020}
UMA Project.
\newblock The Optimistic Oracle.
\newblock UMA Technical Documentation, 2020.

\end{thebibliography}

\end{document}
